\documentclass[aspectratio=169,10pt,english]{beamer}

\input{../../../_preambulo_beamer.tex}
\input{../../../_comandos_beamer.tex}

\selectlanguage{english}

\title{MA1001B - Statistical Modeling for Decision Making}
\subtitle{Lesson 29: Confidence Interval for a Population Proportion}
\author{}
\institute{Tecnol\'ogico de Monterrey}
\date{}

\begin{document}

\begin{frame}[plain]
  \vspace{1.2cm}
  {\Large\bfseries MA1001B - Statistical Modeling for Decision Making\par}
  \vspace{0.7cm}
  {\large Lesson 29: Confidence Interval for a Population Proportion\par}
  \vspace{0.7cm}
  {\color{TecRojo}\rule{\textwidth}{1.2pt}}
  \vspace{0.8cm}

  MA1001B Faculty\\[0.3cm]
  Term\\[0.3cm]
  Tecnol\'ogico de Monterrey
\end{frame}

\begin{frame}{The Binary Rate Question}
  \begin{alertblock}{Guiding question}
    How is a 95\% interval for an unknown proportion built from $x$ successes
    in $n$ independent trials?
  \end{alertblock}

  \vspace{0.6em}
  By the end of this lesson, you can:
  \begin{itemize}
    \item compute $\hat{p}=x/n$;
    \item check $n\hat{p}$ and $n(1-\hat{p})$ before using $z$;
    \item form the Wald interval $\hat{p}\pm 1.96\sqrt{\hat{p}(1-\hat{p})/n}$;
    \item explain why a rare event with small $x$ makes Wald misleading.
  \end{itemize}

  \vfill
  \scriptsize All ticket counts used today are synthetic. No real company data are used.
\end{frame}

\begin{frame}{A Synthetic First-Contact Snapshot}
  \small
  \begin{center}
    \begin{tabular}{lr}
      \toprule
      \textbf{Quantity} & \textbf{Value} \\
      \midrule
      Tickets $n$ & 150 \\
      First-contact resolutions $x$ & 27 \\
      Not first-contact & 123 \\
      Sample proportion $\hat{p}$ & $27/150=0.180000$ \\
      \bottomrule
    \end{tabular}
  \end{center}

  \vspace{0.5em}
  \begin{exampleblock}{Success-failure check}
    $n\hat{p}=27\ge 5$ and $n(1-\hat{p})=123\ge 5$. The normal approximation
    to $\hat{p}$ is treated as usable for this snapshot.
  \end{exampleblock}
\end{frame}

\begin{frame}[fragile]{Step 1: Sample Proportion}
  \begin{columns}[T]
    \begin{column}{0.42\textwidth}
      \small
      \begin{lstlisting}
phat = x / n
n_phat = n * phat
n_qhat = n * (1 - phat)
      \end{lstlisting}

      \begin{block}{Verified output}
        $\hat{p}=0.180000$\\
        $n\hat{p}=27$\\
        $n(1-\hat{p})=123$
      \end{block}
    \end{column}
    \begin{column}{0.58\textwidth}
      \centering
      \includegraphics[width=\textwidth]{../figures/prop_ci_01_sample_proportion.png}
      \vspace{0.2em}
      \scriptsize First-contact counts from Step 1
    \end{column}
  \end{columns}
\end{frame}

\begin{frame}{Wald Interval After the Check Passes}
  \small
  Standard error and 95\% margin:
  \[
    \mathrm{SE}=\sqrt{\frac{0.18\times 0.82}{150}}=0.031369,
  \]
  \[
    E=1.96\times 0.031369=0.061483.
  \]
  Interval:
  \[
    \hat{p}\pm E=[0.118517, 0.241483].
  \]
  The interval estimates the unknown population rate $p$, not the next ticket.
\end{frame}

\begin{frame}[fragile]{Step 2: The 95\% Wald Interval}
  \begin{columns}[T]
    \begin{column}{0.40\textwidth}
      \small
      \begin{lstlisting}
se = np.sqrt(
    phat * (1 - phat) / n
)
lower = phat - 1.96 * se
      \end{lstlisting}

      \begin{block}{Verified output}
        $\mathrm{SE}=0.031369$\\
        Margin $=0.061483$\\
        CI $[0.118517, 0.241483]$
      \end{block}
    \end{column}
    \begin{column}{0.60\textwidth}
      \centering
      \includegraphics[width=\textwidth]{../figures/prop_ci_02_wald_interval.png}
      \vspace{0.2em}
      \scriptsize Wald interval from Step 2
    \end{column}
  \end{columns}
\end{frame}

\begin{frame}{Step 3: Rare Events Make Wald Misleading}
  \begin{columns}[T]
    \begin{column}{0.42\textwidth}
      \small
      Now $x=3$, $n=150$, $\hat{p}=0.020000$.

      \vspace{0.4em}
      $n\hat{p}=3<5$. Conditions fail.

      \vspace{0.4em}
      Wald interval:
      \[
        [-0.002405, 0.042405].
      \]
      The lower bound is negative.

      \vspace{0.5em}
      \textbf{Limit:} a proportion cannot be negative. Small $x$ makes the
      Wald interval misleading.
    \end{column}
    \begin{column}{0.58\textwidth}
      \centering
      \includegraphics[width=\textwidth]{../figures/prop_ci_03_rare_event_wald.png}
      \vspace{0.2em}
      \scriptsize Wald versus Wilson diagnostic from Step 3
    \end{column}
  \end{columns}
\end{frame}

\begin{frame}{Concept Check}
  \begin{enumerate}
    \item Compute $\hat{p}$ from $x=27$ and $n=150$.
    \item Do $n\hat{p}=27$ and $n(1-\hat{p})=123$ pass the cutoff 5?
    \item Why is a negative Wald bound a problem for a proportion?
    \item What quantity fails the check when $x=3$ and $n=150$?
  \end{enumerate}

  \vspace{0.6em}
  \begin{block}{Connection to Lesson 30}
    The session can continue directly into choosing $n$ so that a target
    margin of error is achievable under a budget.
  \end{block}

  \vfill
  \scriptsize Source: OpenStax, \textit{Introductory Business Statistics 2e},
  Chapter 8, Section 8.3. CC BY-NC-SA.
\end{frame}

\appendix



\end{document}
